% ============================================================
% Section 129: 本征半导体态密度与载流子浓度
% ============================================================
\section{半导体物理：抛物带态密度与载流子浓度}
\label{sec:parabolic-dos-carrier-density}

\begin{modelbox}[题目：抛物带近似下的载流子统计]
考虑一本征半导体，设 $E$ 表示一个电子的能量，$g_c(E)$ 表示导带的态密度，$g_v(E)$ 表示价带的态密度。假设 $E_c-E_F\gg k_BT$ 和 $E_F-E_v\gg k_BT$，
\[
g_c(E)=c_1(E-E_c)^{1/2},\qquad g_v(E)=c_2(E_v-E)^{1/2},
\]
其中 $E_c$ 和 $E_v$ 分别表示导带底和价带顶的能量，$E_F$ 是费米能。

(1) 用 $k_B,T,c_1,E_c,E_F$ 和一个无量纲积分表示导带的电子密度；

(2) 用 $k_B,T,c_2,E_v,E_F$ 和一个无量纲积分表示价带的空穴密度；

(3) 求出 $E_F(T)$ 的明确的表达式；

(4) 如果材料用施主原子掺杂，那么在 (1)、(2)、(3) 的结果中有哪一个（如果有的话）仍能成立？解释之。
\end{modelbox}

\begin{tipbox}[考点快速分析]
\begin{itemize}
    \item \textbf{抛物带态密度}：三维 $g(E)\propto(E-E_0)^{1/2}$
    \item \textbf{玻尔兹曼近似}：非简并条件下费米分布退化为指数形式
    \item \textbf{Gamma 函数}：$\int_0^\infty x^{1/2}e^{-x}\,dx=\Gamma(3/2)=\sqrt{\pi}/2$
    \item \textbf{掺杂影响}：破坏 $n=p$ 对称性，费米能级移动
\end{itemize}
\end{tipbox}

\begin{derivationbox}[标准解题过程]
\textbf{(1) 导带电子密度 $n$}

在非简并条件下，$f(E)\approx e^{-(E-E_F)/k_BT}$。电子浓度：
\[
n=\int_{E_c}^\infty c_1(E-E_c)^{1/2}e^{-(E-E_F)/k_BT}\,dE.
\]
令 $x=(E-E_c)/k_BT$，则 $E=E_c+k_BTx$，$dE=k_BT\,dx$：
\[
n=c_1(k_BT)^{3/2}e^{-(E_c-E_F)/k_BT}\int_0^\infty x^{1/2}e^{-x}\,dx.
\]
其中无量纲积分 $\int_0^\infty x^{1/2}e^{-x}\,dx=\Gamma(3/2)=\sqrt{\pi}/2$。因此
\begin{equation}
\boxed{n=\frac{\sqrt{\pi}}{2}\,c_1(k_BT)^{3/2}e^{-(E_c-E_F)/k_BT}.}
\end{equation}

\textbf{(2) 价带空穴密度 $p$}

空穴占据概率 $1-f(E)\approx e^{-(E_F-E)/k_BT}$：
\[
p=\int_{-\infty}^{E_v}c_2(E_v-E)^{1/2}e^{-(E_F-E)/k_BT}\,dE.
\]
令 $x=(E_v-E)/k_BT$，则 $E=E_v-k_BTx$，$dE=-k_BT\,dx$：
\[
p=c_2(k_BT)^{3/2}e^{-(E_F-E_v)/k_BT}\int_0^\infty x^{1/2}e^{-x}\,dx.
\]
因此
\begin{equation}
\boxed{p=\frac{\sqrt{\pi}}{2}\,c_2(k_BT)^{3/2}e^{-(E_F-E_v)/k_BT}.}
\end{equation}

\textbf{(3) 本征费米能级 $E_F(T)$}

本征条件 $n=p$：
\[
c_1e^{-(E_c-E_F)/k_BT}=c_2e^{-(E_F-E_v)/k_BT}.
\]
取对数并整理：
\[
2E_F-E_c-E_v=k_BT\ln\frac{c_1}{c_2}.
\]
解得
\begin{equation}
\boxed{E_F=\frac{E_c+E_v}{2}+\frac{k_BT}{2}\ln\frac{c_1}{c_2}.}
\end{equation}
若 $c_1=c_2$（即电子和空穴有效质量相等），费米能级严格位于禁带中央。一般情况下因 $c\propto(m^*)^{3/2}$，本征费米能级会略微偏离禁带中央。

\textbf{(4) 施主掺杂后的影响}

\textit{(1) 和 (2) 的结果仍成立}：只要半导体处于热平衡且满足非简并条件，导带电子浓度和价带空穴浓度的表达式形式不变，仅 $E_F$ 的具体数值改变。

\textit{(3) 的结果不再成立}：掺入施主后电中性条件变为 $n=p+N_D^+$（$N_D^+$ 为电离施主浓度），$n\neq p$。费米能级将向导带底移动，需由新的电中性方程重新求解。
\end{derivationbox}

\begin{formulabox}[答案汇总]
\begin{center}
\begin{tabular}{ll}
\toprule
\textbf{问题} & \textbf{答案} \\
\midrule
(1) 导带电子密度 & $n=\dfrac{\sqrt{\pi}}{2}\,c_1(k_BT)^{3/2}e^{-(E_c-E_F)/k_BT}$ \\
(2) 价带空穴密度 & $p=\dfrac{\sqrt{\pi}}{2}\,c_2(k_BT)^{3/2}e^{-(E_F-E_v)/k_BT}$ \\
(3) 本征费米能级 & $E_F=\dfrac{E_c+E_v}{2}+\dfrac{k_BT}{2}\ln\dfrac{c_1}{c_2}$ \\
(4) 掺杂影响 & (1)(2) 公式形式仍成立；(3) 不再成立 \\
\bottomrule
\end{tabular}
\end{center}
\end{formulabox}

\begin{insightbox}[物理图像]
\begin{itemize}
    \item 抛物带近似下，载流子浓度正比于 $T^{3/2}e^{-\Delta E/k_BT}$，指数因子主导温度依赖。
    \item 本征费米能级由导带和价带的有效态密度之比决定，通常靠近禁带中央。掺杂破坏电中性对称性，费米能级移动，是半导体导电性可调的根本原因。
    \item 非简并统计公式的形式在掺杂后保持不变，仅 $E_F$ 改变，这大大简化了掺杂半导体的理论处理。
\end{itemize}
\end{insightbox}
